\section*{第 \ref{ch:g7:short-circuits-safety} 章——短路与用电安全}

\begin{solution}{exo:g7:short-circuits-safety:1}
在连接同两点的两条路径之间，电流
更多地涌向阻碍较小的那一条。正在工作的\omterm{def:g3:first-electric-circuit:bulb}{灯泡}会阻碍电流，
而粗裸导线几乎没有阻碍——于是电流离开\omterm{def:g3:first-electric-circuit:bulb}{灯泡}，
转向导线。
\end{solution}

\begin{solution}{exo:g7:short-circuits-safety:2}
器件\omterm{def:g7:short-circuits-safety:device}{短路}：在某个器件两端并接一条旁路，该器件停止工作，
而回路仍然连通。\omterm{def:g3:first-electric-circuit:battery}{电池}\omterm{def:g7:short-circuits-safety:device}{短路}：一条裸导线连接了
\omterm{def:g3:first-electric-circuit:battery}{电池}自身的两个\omterm{def:g3:first-electric-circuit:battery}{接线柱}，中间不接任何器件——这是更严重的情况：
电路中无任何阻碍，大电流猛冲，使导线与
\omterm{def:g3:first-electric-circuit:battery}{电池}急剧发热，最终可能损坏甚至起火。
\end{solution}

\begin{solution}{exo:g7:short-circuits-safety:3}
旁路把\omterm{def:g3:first-electric-circuit:bulb}{灯泡} 2 的全部电流分走，
\omterm{def:g3:first-electric-circuit:bulb}{灯泡}随之熄灭。由于其中一个阻碍被旁路，整个回路对\omterm{def:g3:first-electric-circuit:battery}{电池}的阻碍减小，
电流随之增强，\omterm{def:g3:first-electric-circuit:bulb}{灯泡} 1 独自承受更大的电流，
因而变得更亮。
\end{solution}

\begin{solution}{exo:g7:short-circuits-safety:4}
电流通过导线时总会使导线发热，大电流发热
尤为剧烈。\omterm{def:g3:first-electric-circuit:bulb}{灯泡}正是有意利用这种发热：灯丝
被设计成能够炽热发白而发光。
\end{solution}

\begin{solution}{exo:g7:short-circuits-safety:5}
一根特意做得很细的牺牲导线——\omterm{def:g6:water-states:changes}{熔化}而切断
回路，当电流超过额定值时，它便自我牺牲。
\omterm{def:g7:short-circuits-safety:fuse}{保险丝}应\omterm{def:g4:series-circuits:series}{串联}接入，因为大电流必须穿过它，
\omterm{def:g5:series-parallel-circuits:parallel}{并联}的\omterm{def:g7:short-circuits-safety:fuse}{保险丝}只能眼看着电流从旁边溜过。
\end{solution}

\begin{solution}{exo:g7:short-circuits-safety:6}
断路器不靠\omterm{def:g6:water-states:changes}{熔化}，而是跳闸断开，
复位时只需拨动一下，无需更换零件。在复位或更换之前，
必须先找出并排除故障：保护装置跳闸是有原因的。
\end{solution}

\begin{solution}{exo:g7:short-circuits-safety:7}
过载：三件无辜的用电器同时要求同一条\omterm{def:g5:series-parallel-circuits:parallel}{支路}
承载它们各自满负荷的三股电流，总电流使导线
发热——断路器按总电流跳闸，这是正确的。
\end{solution}

\begin{solution}{exo:g7:short-circuits-safety:8}
钉子永远不会\omterm{def:g6:water-states:changes}{熔化}：它使保护装置失效，而隐患
仍在蔓延。下一次大电流遇不到任何关卡，房屋的
墙内布线本身充当了\omterm{def:g7:short-circuits-safety:fuse}{保险丝}——在无人察觉处
悄悄烧焦。
\end{solution}

\begin{solution}{exo:g7:short-circuits-safety:9}
器件上游发生\omterm{def:g7:short-circuits-safety:device}{短路}——从该电线看去，相当于电源被短接，
\omterm{def:g3:first-electric-circuit:bulb}{灯泡}熄灭，大电流穿过磨损的接触点，
使电线发热，
而该\omterm{def:g5:series-parallel-circuits:parallel}{支路}的\omterm{def:g7:short-circuits-safety:fuse}{保险丝}或断路器随即
切断该\omterm{def:g5:series-parallel-circuits:parallel}{支路}。这正说明系统在正常工作。
\end{solution}

\begin{solution}{exo:g7:short-circuits-safety:10}
这根导线横跨两只\omterm{def:g3:first-electric-circuit:bulb}{灯泡}，经普通导线和\omterm{def:g3:first-electric-circuit:battery}{电池}引线，
将\omterm{def:g3:first-electric-circuit:battery}{电池}自身的两个\omterm{def:g3:first-electric-circuit:battery}{接线柱}相连——
旁路上不接任何器件：这是\omterm{def:g3:first-electric-circuit:battery}{电池}\omterm{def:g7:short-circuits-safety:device}{短路}。
两只\omterm{def:g3:first-electric-circuit:bulb}{灯泡}都会熄灭，大电流使导线和\omterm{def:g3:first-electric-circuit:battery}{电池}发热；即便是
在玩具规模下，它也会损坏\omterm{def:g3:first-electric-circuit:battery}{电池}、烫手，
因此只能以文字说明，不宜亲手演示。
\end{solution}

\begin{solution}{exo:g7:short-circuits-safety:11}
熄灭\omterm{def:g3:first-electric-circuit:bulb}{灯泡}的支撑丝已熔接在一起，成了内置的旁路：
回路仍保持闭合，电流\emph{经由}熄灭\omterm{def:g3:first-electric-circuit:bulb}{灯泡}的\omterm{def:g7:short-circuits-safety:device}{短路}，
\omterm{def:g4:series-circuits:series}{串联电路}因此不致全部断开——少了一个阻碍，所以
略微变亮。但每次这样的熄灭，都会增强流经
其余\omterm{def:g3:first-electric-circuit:bulb}{灯泡}的电流，从而加速下一次损坏：这串\omterm{def:g3:first-electric-circuit:bulb}{灯泡}
悄然酝酿自身的连锁失效。
\end{solution}

\begin{solution}{exo:g7:short-circuits-safety:12}
干路上设一个总保护装置和总\omterm{ex:g3:first-electric-circuit:switch}{开关}；随后每条\omterm{def:g5:series-parallel-circuits:parallel}{支路}
——照明、工具、充电器——各自带有独立\omterm{ex:g3:first-electric-circuit:switch}{开关}，
以及与其容量相匹配的保护装置（工具\omterm{def:g5:series-parallel-circuits:parallel}{支路}承载的
电流最大，应配最坚固的断路器；充电器的
细\omterm{def:g5:series-parallel-circuits:parallel}{支路}需要最灵敏的\omterm{def:g7:short-circuits-safety:fuse}{保险丝}，因为它会在
总保护装置忽略的电流下烧焦）。例如在抽屉上方贴一张提示：
“熔断的\omterm{def:g7:short-circuits-safety:fuse}{保险丝}说的是实话——查找故障所在。
同规格更换；锡箔和铁钉什么也保护不了，只会
为火灾开路。”
\end{solution}

\begin{solution}{pb:g7:short-circuits-safety:1}
\textbf{1.} 过载：三件耗能大的用电器把各自正常的电流
叠加在同一条\omterm{def:g5:series-parallel-circuits:parallel}{支路}上——无需裸导线，单是这股总电流
就足以使线路发热到超过安全限度。
\textbf{2.} 咬破处发生\omterm{def:g7:short-circuits-safety:device}{短路}：从那一刻起，电流
离开了其后方的\omterm{def:g3:first-electric-circuit:bulb}{灯泡}，大电流穿过相互接触的
\omterm{def:g4:conductors-insulators:conductor}{导体}，使电线发热。
\textbf{3.} 它取消了保护：锡箔能导通任何大电流，
且不熔断——该\omterm{def:g5:series-parallel-circuits:parallel}{支路}于是失去保护。
\textbf{4.} 大致经过：先是过载，真正的\omterm{def:g7:short-circuits-safety:fuse}{保险丝}
熔断了——这是正常的——也许不止一次；有人
用锡箔“解决”了这个麻烦（见第三件证物）；等
被啃坏的电线最终\omterm{def:g7:short-circuits-safety:device}{短路}时，已无任何保护，于是
大电流把墙内线路烧成了火，而不是在黑暗中
随\omterm{def:g7:short-circuits-safety:fuse}{保险丝}的熔断而戛然而止。
\textbf{5.} 电流总会使其所经的导线发热；冲得越猛，
发热越厉害。\omterm{def:g3:first-electric-circuit:bulb}{灯泡}证明了这种发热是可以驯服的——其灯丝
在安全的\omterm{def:g3:first-electric-circuit:bulb}{灯泡}中被设计为白炽发光。这里缺少的是
让\omterm{def:g5:heat-and-insulation:heat}{热量}有所安放的设计，以及阻止它的任何保护装置。
\textbf{6.} 特意做细，使它在电路所有导线中
最先\omterm{def:g6:water-states:changes}{熔化}：保护装置把致命的发热集中在自己
三\omterm{def:g2:measuring-length:units}{厘米}长的那一段中——以自己的牺牲，使墙内的数米导线不致受损。
\textbf{7.} 保护装置必须在恰当的时刻\emph{失效}：它的
优点恰恰在于经过设计的薄弱环节。锡箔作为
\omterm{def:g4:conductors-insulators:conductor}{导体}的优良性能，正是对保护关卡的败坏——
通行无阻，却毫无保护。
\textbf{8.} 安全并不像礼节那样可以相加：三股合乎规定的
电流在同一条路上相加，便成了一股违规的大电流——\omterm{def:g5:series-parallel-circuits:parallel}{支路}
服从的是算术，而非人的意图。
\textbf{9.} 锡箔的问题。过载和被啃坏的电线都是
系统设计时就应能承受的意外——保护装置正是为此而设；
让保护装置失效，把本可承受的故障变成了
灾难。
\textbf{10.} 若安装真正的\omterm{def:g7:short-circuits-safety:fuse}{保险丝}：电线\omterm{def:g7:short-circuits-safety:device}{短路}，大电流
骤然涌来，\omterm{def:g7:short-circuits-safety:fuse}{保险丝}瞬间熔断——厨房陷入
黑暗；家人回来后，看到的是熔断的\omterm{def:g7:short-circuits-safety:fuse}{保险丝}和被啃坏的电线，
嘀咕几句，把两者修好，却从不知道当晚的情况
险些有多严重。
\textbf{11.} 每个发热电器各插在自己的墙上插座上，接在各自
保护足够的\omterm{def:g5:series-parallel-circuits:parallel}{支路}上；插线板仅用于
灯具和充电器，并遵守一条长期规定：每个插座只接
一件耗电大的用电器，插线板绝不可给电暖器供电。
\textbf{12.} 例如：“坏人的手法很简单：
为电流提供一条旁路，它便会走此捷径，
并在冲过时发热。保护装置的职责，是\omterm{def:g4:series-circuits:series}{串联}在电路中，
率先、坦然而低成本地失效。任何家庭都不得
随意自制保护装置：\omterm{def:g7:short-circuits-safety:fuse}{保险丝}须以\omterm{def:g7:short-circuits-safety:fuse}{保险丝}更换——
同规格更换，否则墙壁终有一天会自己\omterm{def:g6:water-states:changes}{熔化}。”
\end{solution}
